% Exercise 2.14
If we compare the two methods of depreciation over a cycle of three years, we have
for $\text{AV}$ the initial value of the asset and $\text{PV}_{\text{A}}$ and
$\text{PV}_{\text{B}}$ the present values of depreciation after three years, that
\begin{equation*}
    \begin{split}
        \text{PV}_{\text{A}} &= \text{AV} \sum_{k=1}^3{\frac{d_k}{({1 + r})^k}} \\
                             &= \text{AV} \Biggl[\frac{.25}{1 + r} 
                                 + \frac{.38}{{(1 + r)}^2} 
                                 + \frac{.37}{{(1 + r)}^3}\Biggr].
    \end{split}
\end{equation*}
and
\begin{equation*}
    \begin{split}
        \text{PV}_{\text{B}} &= \text{AV} \sum_{k=1}^3{\frac{d_k}{({1 + r})^k}} \\
                             &= \text{AV} \sum_{k=1}^3{\frac{0.\dot{3}}{({1 + r})^k}},
    \end{split}
\end{equation*}
for $r$ the assumed rate of interest, so that the method which maximizes the 
present value of depreciation depends largely on the interest rate, since
this will determine the relative superiority in present value of depreciation.
